\documentclass[../principal.tex]{subfiles}
\graphicspath{{\subfix{../imagenes/}}}

\begin{document}

\ifSubfilesClassLoaded{
   \chapter{Circuitos integrados}
}{}

\section{Chip 4093}

El 4093 contiene 4 compuertas lógicas NAND de 2 entradas con activador de Schmitt.

Lo usaremos para hacer osciladores de onda cuadrada.

\begin{figure}[H]
\centering
\begin{circuitikz}
  % cuerpo del chip
  \node[draw, thick, minimum width=3cm, minimum height=8cm] (chip) at (0,0) {\rotatebox{-90}{\large 4093}};
  
  % patitas izquierda
  \foreach \y/\text in {
    3.0/{01 $A_{1}$},
    2.0/{02 $B_{1}$},
    1.0/{03 $Q_{1}$},
    0.0/{04 $Q_{2}$},
   -1.0/{05 $A_{2}$},
   -2.0/{06 $B_{2}$},
   -3.0/{07 GND},
  }{
    % patita
    \draw ($(chip.west)+(0,\y)$) -- ++(-0.4,0);
    % etiqueta
    \node[left] at ($(chip.west)+(-0.4,\y)$) {\text};
  }

  % patitas derecha
  \foreach \y/\text in {
    3.0/{14 $V_{CC}$},
    2.0/{13 $B_{4}$},
    1.0/{12 $A_{4}$},
    0.0/{11 $Q_{4}$},
   -1.0/{10 $Q_{3}$},
   -2.0/{09 $B_{3}$},
   -3.0/{08 $A_{3}$},
  }{
    % patita
    \draw ($(chip.east)+(0,\y)$) -- ++(0.4,0);
    % etiqueta
    \node[right] at ($(chip.east)+(0.4,\y)$) {\text};
  }

  % puntito en patita 1
  \fill ($(chip.north west)+(0.3,-0.3)$) circle (2pt);
\end{circuitikz}
\caption{Integrado 4093}
\end{figure}

\subsection{Patitas}

El chip 4093 tiene 14 patitas.

Dos patitas son para alimentación: 14 para $V_{CC}$ y 07 para GND.

Las otras patitas corresponden a entradas y salidas de las 4 compuertas NAND, donde cada compuerta tiene 2 entradas (A y B) y una salida (Q), y están numeradas de 1 a 4.

\begin{itemize}
    \item 01: $A_{1}$
    \item 02: $B_{1}$
    \item 03: $Q_{1}$
    \item 04: $Q_{2}$
    \item 05: $A_{2}$
    \item 06: $B_{2}$
    \item 07: GND
    \item 08: $A_{3}$
    \item 09: $B_{3}$
    \item 10: $Q_{3}$
    \item 11: $Q_{4}$
    \item 12: $A_{4}$
    \item 13: $B_{4}$
    \item 14: $V_{CC}$
\end{itemize}

\subsection{Osciladores de onda cuadrada}

Para usar el 4093 como osciladores, conectaremos la alimentación del circuito integrado, y luego cada compuerta de un modo que sea un oscilador de onda cuadrada.

La tabla de verdad de cada compuerta NAND es la siguiente:

\begin{center}
\begin{tabular}{cc|c}
$A$ & $B$ & $Q$ \\
\hline
0 & 0 & 1 \\
0 & 1 & 1 \\
1 & 0 & 1 \\
1 & 1 & 0 \\
\end{tabular}
\end{center}

Si nos concentramos en la entrada A, vemos que cuando es 0, la salida Q siempre vale 1, lo que no nos sirve para nuestra intención de oscilar.

En cambio, cuando la entrada A es 1, la salida Q es el inverso de la entrada B, lo que sí permite la oscilación. 

Usaremos entonces una entrada A fija en $V_{CC}$ cuando queremos que la compuerta oscile siempre, o una entrada A conectada a un secuenciador que permita la oscilación solamente cuando esta entrada sea $V_{CC}$.


\subsubsection{Alimentación}

\begin{enumerate}
    \item Patita 14 a $V_{CC}$
    \item Patita 07 a GND
\end{enumerate}

\subsubsection{Entradas}

\begin{enumerate}
    \item Patita 01 a $V_{CC}$ o a un secuenciador que controle compuerta 1
    \item Patita 05 a $V_{CC}$ o a un secuenciador que controle compuerta 2
    \item Patita 08 a $V_{CC}$ o a un secuenciador que controle compuerta 3
    \item Patita 12 a $V_{CC}$ o a un secuenciador que controle compuerta 4
\end{enumerate}

\subsubsection{Salidas}

\begin{enumerate}
    \item Patita 03 - salida de la compuerta 1
    \item Patita 04 - salida de la compuerta 2
    \item Patita 10 - salida de la compuerta 3
    \item Patita 11 - salida de la compuerta 4
\end{enumerate}

\subsubsection{Retroalimentación}

Conectaremos las entradas B de cada compuerta a su salida Q correspondiente, a través de un circuito RC.

Para lograrlo, cada entrada B se conecta a tierra a través de un capacitor, y también se conecta a su salida Q a través de un potenciómetro.

Los valores de los capacitores y los potenciómetros son los que determinan la frecuencia de oscilación de cada compuerta NAND.

\begin{enumerate}
    \item Patita 02 conectada a tierra a través de un capacitor de 1 $\mu F$
    \item Patita 06 conectada a tierra a través de un capacitor de 1 $\mu F$
    \item Patita 09 conectada a tierra a través de un capacitor de 1 $\mu F$
    \item Patita 13 conectada a tierra a través de un capacitor de 1 $\mu F$
    \item Patita 02 conectada a patita 03 a través de un potenciómetro de 100 k$\Omega$
    \item Patita 06 conectada a patita 04 a través de un potenciómetro de 100 k$\Omega$
    \item Patita 09 conectada a patita 10 a través de un potenciómetro de 100 k$\Omega$
    \item Patita 13 conectada a patita 11 a través de un potenciómetro de 100 k$\Omega$
\end{enumerate}

\newpage

\end{document}